\documentclass[11pt]{amsart}

\usepackage{amsmath, amssymb, amsthm}
\usepackage{mathrsfs}
\usepackage{hyperref}
\usepackage{cleveref}

%---------- THEOREM ENVIRONMENTS ----------
\newtheorem{theorem}{Theorem}[section]
\newtheorem{proposition}[theorem]{Proposition}
\newtheorem{lemma}[theorem]{Lemma}
\newtheorem{corollary}[theorem]{Corollary}

\theoremstyle{definition}
\newtheorem{definition}[theorem]{Definition}
\newtheorem{remark}[theorem]{Remark}
\newtheorem{example}[theorem]{Example}

%---------- TITLE ----------
\title{On the Structure and Compactness Properties of Asymptotic $L^p$ Spaces}

\author{Hikmatullo Ismatov}
\address{King Abdullah University of Science and Technology (KAUST)}
\email{your@email.com}

%---------- ABSTRACT ----------
\begin{document}

\begin{abstract}
We study asymptotic $L^p$ spaces, recently introduced by Alves, which provide a natural framework for analyzing measurable functions on infinite measure spaces. These spaces interpolate between classical $L^p$ spaces and convergence in measure, and are equipped with a nonlocally convex $F$-norm structure.

In this work, we investigate structural and compactness properties of asymptotic $L^p$ spaces. We establish embedding results, develop compactness criteria extending the Kolmogorov--Riesz theorem, and provide new formulations in terms of almost equicontinuity and tail control. Our results contribute to the development of a function space theory adapted to truncated $L^p$ structures.
\end{abstract}

\maketitle

%---------- INTRODUCTION ----------
\section{Introduction}

Classical $L^p$ spaces play a central role in analysis; however, on infinite measure spaces they often impose overly restrictive conditions, while convergence in measure lacks sufficient structure for many analytical purposes. To address this gap, Alves introduced the notion of asymptotic $L^p$ convergence, which combines features of both $L^p$ convergence and convergence in measure.

More precisely, this convergence can be characterized by the truncated functional
\[
\int_X \min(|f|,1)^p \, d\mu,
\]
which induces an $F$-norm topology on the space of measurable functions. This framework has been further developed in recent works, where compactness criteria and structural properties have been established.

The aim of this paper is to further develop the functional analytic theory of asymptotic $L^p$ spaces. Our main contributions are as follows:

\begin{itemize}
\item We establish embedding results connecting Sobolev-type spaces and asymptotic $L^p$ spaces.
\item We prove compactness criteria extending the Kolmogorov--Riesz theorem in this setting.
\item We provide alternative formulations of compactness in terms of almost equicontinuity and tail control.
\end{itemize}

These results contribute to the understanding of asymptotic $L^p$ spaces as a natural intermediate framework between $L^p$ theory and measure-theoretic convergence.

%---------- PRELIMINARIES ----------
\section{Preliminaries}

Let $(X,\mathcal{A},\mu)$ be a measure space.

\begin{definition}
For $1 \le p < \infty$, the asymptotic $L^p$ functional is defined by
\[
\|f\|_{\alpha_p}
:=
\left(
\int_X \min(|f(x)|,1)^p \, d\mu(x)
\right)^{1/p}.
\]
\end{definition}

\begin{definition}
A sequence $(f_n)$ converges asymptotically in $L^p$ to $f$ if
\[
\|f_n - f\|_{\alpha_p} \to 0.
\]
\end{definition}

We recall the following characterization.

\begin{proposition}
A sequence $(f_n)$ converges asymptotically in $L^p$ to $f$ if and only if
\[
\int_X \min(|f_n - f|,1)^p \, d\mu \to 0.
\]
\end{proposition}

%---------- BASIC PROPERTIES ----------
\section{Basic Properties of the $\alpha_p$-Functional}

\begin{proposition}
The functional $\|\cdot\|_{\alpha_p}$ defines an $F$-norm on the space of measurable functions.
\end{proposition}

\begin{proof}
[Sketch] The proof follows from the truncation structure of the integrand and standard properties of measurable functions.
\end{proof}

%---------- EMBEDDING ----------
\section{Embedding Results}

\begin{theorem}
Let $u \in L^p(\mathbb{R}^n)$. Then $u \in \Lambda_p(\mathbb{R}^n)$ and
\[
\|u\|_{\alpha_p} \le \|u\|_{L^p}.
\]
\end{theorem}

%---------- COMPACTNESS ----------
\section{Compactness in Asymptotic $L^p$ Spaces}

\begin{definition}
A family $\mathcal{F}$ is almost equibounded if for every $\varepsilon>0$ there exists $M>0$ such that
\[
\sup_{f \in \mathcal{F}} \mu(\{|f|>M\}) < \varepsilon.
\]
\end{definition}

\begin{theorem}
Let $\mathcal{F} \subset \Lambda_p(\mathbb{R}^n)$. Suppose:
\begin{enumerate}
\item $\mathcal{F}$ is almost equibounded,
\item $\mathcal{F}$ satisfies a uniform tail condition,
\item $\mathcal{F}$ satisfies a translation continuity condition.
\end{enumerate}
Then $\mathcal{F}$ is relatively compact in $\Lambda_p$.
\end{theorem}

%---------- FURTHER RESULTS ----------
\section{Further Results}

We discuss refinements of the compactness criterion and possible formulations in terms of almost equicontinuity.

%---------- REFERENCES ----------
\begin{thebibliography}{9}

\bibitem{Alves2024}
N. J. Alves,
\emph{On $F$-spaces of almost-Lebesgue functions},
arXiv:2410.22905.

\bibitem{AlvesOniani2024}
N. J. Alves and E. Oniani,
\emph{Asymptotic $L_p$ convergence and related topics},
arXiv:2407.06830.

\bibitem{Alves2025}
N. J. Alves,
\emph{A Kolmogorov--Riesz compactness theorem in asymptotic $L_p(\mathbb{R}^n)$},
arXiv:2507.15102.

\end{thebibliography}

\end{document}
